\section{模型检验与结果分析}
\subsection{指标与证据来源}
表\ref{tab:demo}是根据解析公式计算的教学示例，不含官方测试、运行时间或真实任务成功率。
正式实验需写数据来源、版本、样本规模、统计分母、指标公式和计时起止点。
开发集反复用于调参后不能继续充当未见独立确认集。

\begin{table}[htbp]
\centering
\caption{Teaching example: exact objective values}
\label{tab:demo}
\begin{tabular}{lrrr}
\toprule Method & $x$ & $J(x)$ & $J_{\mathrm{rob}}(x)$ \\
\midrule
baseline & 1.0 & 1.00 & 1.21 \\
analytic & 2.0 & 0.00 & 0.01 \\
\bottomrule
\end{tabular}
\end{table}


\subsection{敏感性、鲁棒性与局限}
本例若将扰动半径替换为 $\delta\geq0$，稳健最优值变为 $\delta^2$，
最优点仍为 $2$。实际问题应报告尾部风险、失败案例与重要分层，不能只写平均提升。
\begin{figure}[htbp]
\centering\placeholderfigure
\caption{图框演示，需替换为真实图并注明数据口径}
\end{figure}
